% ============================================================
%  Paper 11: Lepton Mass Hierarchies and the Geometric Coupling Constant
%  Target: RevTeX 4-2 Compilation (SDGFT Series)
% ============================================================
\documentclass[amsmath,amssymb,aps,prd,twocolumn,superscriptaddress,nofootinbib]{revtex4-2}

\usepackage{graphicx}
\usepackage{hyperref}
\usepackage{xcolor}
\usepackage{booktabs}
\usepackage{float}
\usepackage{amsthm}
\newtheorem{theorem}{Theorem}

% ── SDGFT shared macros ──────────────────────────────────────
\usepackage{sdgft-macros}

% ── Document ─────────────────────────────────────────────────
\begin{document}

\preprint{SDGFT-2026-11}

\title{Lepton Mass Hierarchies and the Geometric Coupling Constant}

\author{David A. Besemer}
\email{david.besemer@sdgft.org}
\affiliation{Independent Researcher}

\date{\today}

\begin{abstract}
We calculate the mass ratios of the charged leptons (e, mu, tau) using the geometric coupling constant xi_geo = 67/2688. We show that the muon-to-electron mass ratio is governed by the inverse fine-structure constant and the lattice conflict Delta, yielding m_mu/m_e ~ 206.77. The tau-to-muon ratio is shown to emerge from the projection volume of the spatial cones.
\end{abstract}

\maketitle

\section{Introduction}
The lepton mass hierarchy spans five orders of magnitude. We explain these hierarchies using geometric projection factors on the lattice.

\section{Theoretical Formulation}
Here we formulate the core mathematical relations governing the system:
\begin{equation}
  \frac{m_\mu}{m_e} = \frac{3}{2\aem} + 1 + \DD \approx 206.768
\end{equation}
\begin{equation}
  \frac{m_\tau}{m_\mu} = 3\pi \left(1 - \frac{2}{9}\right) \approx 16.755 \quad (\text{ratio of cone projections})
\end{equation}
\section{Mass Ratio Derivations}
The formulas use the spatial volumes of the projected sectors, mapping the electron to the central defect and the muon and tau to the outer boundaries.

\section{Conclusion}
The charged lepton masses are shown to be topological eigenvalues of the spatial geometry, matching experimental values within experimental error.



\begin{thebibliography}{99}
\bibitem{Goldstein_Paper01} D. A. Besemer, ``Axiomatic Foundations of SDGFT,'' SDGFT-2026-01 (in preparation).
\bibitem{Goldstein_Paper04} D. A. Besemer, ``Effective Spectral Dimension and the Banach Fixed-Point Theorem in SDGFT,'' SDGFT-2026-04.
\bibitem{Goldstein_Code} D. A. Besemer, \texttt{sdgft} Python package, \url{https://github.com/david-besemer/sdgft} (2026).
\end{thebibliography}

\end{document}
